\begin{table}[!htbp]
\centering
\caption{Heterogeneity in the association between bullying and reading achievement}
\label{tab:bullying-interactions}
\scriptsize
\setlength{\tabcolsep}{2.2pt}
\resizebox{\textwidth}{!}{%
\begin{tabular}{lcccc}
\toprule
 & \multicolumn{2}{c}{South Africa} & \multicolumn{2}{c}{Brazil} \\
\cmidrule(lr){2-3}\cmidrule(lr){4-5}
 & I1 & I2 & I1 & I2 \\
\midrule
Bullied about monthly & \shortstack{-46.5***\\(7.0)} & \shortstack{-42.2***\\(5.7)} & \shortstack{-34.1***\\(8.8)} & \shortstack{-39.8***\\(5.6)} \\
Bullied about weekly & \shortstack{-95.3***\\(7.8)} & \shortstack{-95.5***\\(6.4)} & \shortstack{-120.4***\\(10.6)} & \shortstack{-110.9***\\(6.9)} \\
Boy & \shortstack{-41.2***\\(6.7)} & \shortstack{-42.9***\\(3.4)} & \shortstack{-4.1\\(4.9)} & \shortstack{-5.6\\(4.7)} \\
Age (years) & \shortstack{-7.6**\\(3.4)} & \shortstack{-8.0**\\(3.4)} & \shortstack{-12.9***\\(4.4)} & \shortstack{-13.5***\\(4.3)} \\
Home SES (ASBHSES scale points) & \shortstack{25.7***\\(2.3)} &  & \shortstack{26.2***\\(2.1)} &  \\
Bullied about monthly × Boy & \shortstack{-0.7\\(9.4)} &  & \shortstack{-10.5\\(11.1)} &  \\
Bullied about weekly × Boy & \shortstack{-5.1\\(7.4)} &  & \shortstack{9.9\\(12.7)} &  \\
Home SES (centred scale points) &  & \shortstack{37.5***\\(2.8)} &  & \shortstack{24.1***\\(2.0)} \\
Bullied about monthly × Home SES (centred scale points) &  & \shortstack{-13.2***\\(2.8)} &  & \shortstack{3.2\\(3.0)} \\
Bullied about weekly × Home SES (centred scale points) &  & \shortstack{-20.3***\\(2.8)} &  & \shortstack{9.0**\\(4.5)} \\
Other gender category &  &  &  & \shortstack{-86.7**\\(35.2)} \\
\midrule
Learners & 9,167 & 9,167 & 4,076 & 4,110 \\
$R^{2}$ & 0.278 & 0.289 & 0.314 & 0.319 \\
\bottomrule
\end{tabular}%
}
\parbox{0.98\textwidth}{\footnotesize \textit{Notes:} Entries are coefficients with jackknife standard errors in parentheses. I1 interacts bullying with gender and excludes Brazil's 41 learners in the Other gender category. I2 interacts bullying with continuous ASBHSES centred at the design-weighted country mean; one unit is one PIRLS ASBHSES scale point. The reference bullying category is almost never and the gender reference is girls. Significance: *** p\textless{}0.01, ** p\textless{}0.05, * p\textless{}0.10.}
\end{table}
